% =============================================================================
% 4. CONCLUSIONES  (15 %)
% Espacio: ~0,5 paginas.
% Que hemos aprendido de la dificultad de los problemas, como escala el coste y limitaciones.
% Material: docs/dificultades_encontradas.md
% =============================================================================
\section{Conclusiones}
\label{sec:conclusiones}

\pendiente{GGH: con una base buena el problema es fácil, y LLL la consigue en las tallas del trabajo.}

\pendiente{LWE: el error es un vector corto; el ataque primal lo encuentra mientras la dimensión
  sea pequeña.}

\pendiente{Ring-LWE: los parámetros de las tareas permiten un atajo (sin reducción mod $q$).}

\pendiente{limitaciones: métodos aproximados (Babai) frente a exactos (enumeración); cómo crece el
  tiempo con la talla; qué haría falta con parámetros reales (dimensiones de cientos).}

\pendiente{dificultades encontradas y decisiones (de \texttt{docs/dificultades\_encontradas.md}),
  por ejemplo LLL demasiado lento y cómo se resolvió.}
